\section{模型假设}

\begin{enumerate}
\item 题目附件中给出的任务需求、设备容量、功率参数、电价、碳强度和新能源出力等数据，在其给定精度内视为可靠的确定性参数；除情景分析中主动设置的扰动外，不考虑设备故障、通信中断等突发事件。
\item 调度以小时为决策步长，任务在整数小时选择开工时刻；任务一旦开始便连续执行，不允许拆分、抢占、暂停或中途迁移，跨小时能耗按照任务与各小时的实际重叠时长计算。
\item 同类任务的单位 GPU 功率、各区域 PUE 和区域间单向网络时延在研究期内保持不变；不考虑网络拥塞、带宽限制以及数据传输产生的额外能耗和费用。
\item 购售电价格、电网碳强度、可用新能源和固定负荷在单个小时内保持恒定，各区域独立满足能源平衡；不考虑跨区域电力潮流、线路损耗及设备全生命周期碳排放。
\item 储能充放电效率采用附件给定值并在研究期内保持不变，忽略储能自放电和容量衰减；为避免通过透支期末电量获得虚假收益，规定调度结束时储能 SOC 恢复至初始水平。
\end{enumerate}
